
\needspace{3\baselineskip}
\color{uniblue}{\subsection{Логика отключения выключателя стороны ВН трансформатора (ЛО ВН)}}
\color{black}

\begin{enumerate}

\item
Логика отключения выключателя стороны ВН трансформатора служит для формирования команды отключения выключателя стороны ВН трансформатора от защит в составе устройства и внешних защит, которые действуют через устройство.

\item
В устройстве логика выключателя стороны ВН трансформатора представлена функциональным блоком «ЛО ВН». Функциональный блок «ЛО ВН» состоит из функции логики отключения («ЛО»).

\item
Ввод функции в работу на этапе параметрирования устройства осуществляется программным переключателем SGF1 «Ввод\_функции». Информация о введенной функции отображается выходным сигналом «ЛО ВН / ЛО: Ввод».

\item
Функция может быть выведена из работы оперативно путем активации сигнала «ОВ ЛО ВН», что характеризуется активным состоянием сигнала «ЛО ВН / ЛО: Оперативный вывод» на выходе функции.

\item
При появлении сигналов срабатывания от общей логики отключения трансформатора и внешних устройств РЗА сторон НН функция формирует сигналы «Отключение» и «Отключить аварийно», сигнал «Отключить аварийно» является импульсным, время импульса определяется выдержкой времени таймера T1 (уставка «Тимп»).

\item
Алгоритм работы функции выполнен в соответствии с рисунком \ref{lovn:img1}. В таблице \ref{lovn:tbl1} приведены параметры, необходимые для настройки функции.

\vspace{2mm}
\begin{figure}[!h]
\centering
% ЗДЕСЬ ВЕТВЛЕНИЕ, в зависимости от типа вставляем разные рисунки
\ifboolexpr{bool {isT}}
{
\includegraphics[width=0.85\textwidth,height=0.85\textheight,keepaspectratio]{\fbpath/15. ЛО В ВН Т 35 ДЗТ/_latex/img/t.pdf}
}{}
\ifboolexpr{bool {isDZT2}}
{
\includegraphics[width=0.85\textwidth,height=0.85\textheight,keepaspectratio]{\fbpath/15. ЛО В ВН Т 35 ДЗТ/_latex/img/dzt.pdf}
}{}
\ifboolexpr{bool {isT2}}
{
\includegraphics[width=0.85\textwidth,height=0.85\textheight,keepaspectratio]{\fbpath/15. ЛО В ВН Т 35 ДЗТ/_latex/img/dzt.pdf}
}{}

%%%%%%%%%%%%%%%%%%%%%%%%%%%%%%%%%%%%%%%%%%%%%%%%%%%%%%%%%%%%%%%%%%%
\captionof{figure}{Функциональная схема ЛО ВН}\label{lovn:img1}
\end{figure}

\end{enumerate}
\FloatBarrier % Все float'ы ДО этой строки будут размещены здесь.

